\section{Local complex waves below the envelope}\label{sec:packets}

Fix a smooth positive input pair and an envelope from
Proposition~\ref{prop:envelope}. A wave packet will have amplitude
comparable to \(e^{kh}\) at its center, decay rapidly relative to
that envelope away from the center, and have a residual as small
as any prescribed inverse power of \(k\). The finite construction
must retain the same derivative and spanning exponents when its
accuracy is increased.

\subsection{Characteristic Hessians and polynomial equations}

In \(G\)-normal coordinates at a center, let
\(p=\nabla h\), \(Q=\sqrt{1+|p|^2}\) and \(H=\nabla^2h\).
For a unit \(v\perp p\), the bilinear square of \(P=p+iQv\) is
\(-1\). Require the phase to satisfy the bilinear eikonal equation
\(G^{-1}(\dd\phi,\dd\phi)=-1\) to first order at the center.
A phase with differential \(P\) must then have Hessian \(S\)
satisfying \(SP=0\): differentiating the equation gives this condition
because the first metric derivatives vanish at the center.
The real Hessian gap \(H-\operatorname{Re}S>0\)
will keep its exponential below \(e^{kh}\).

\begin{lemma}\label{lem:hessian}
Let \(e_1,v\) be orthonormal, \(0\le b<1\), and \(H\) real
symmetric, with
\[
 |H|\le k^{4\delta},\qquad
 b^2H(e_1,e_1)+H(v,v)\ge k^{-2\delta}.
\]
For all sufficiently large \(k\), uniformly under these hypotheses,
there are real symmetric \(U,D\) with
\begin{equation}\label{eq:hessian}
 (U+iD)(be_1+iv)=0,\quad
 H-U\ge k^{-3\delta}I,\quad |U+iD|\le k^{20\delta}.
\end{equation}
\end{lemma}

\begin{proof}
For \(b>0\) the first identity is
\(Dv=bUe_1,\ De_1=-Uv/b\). The prescribed columns are compatible
with symmetry exactly when \(b^2U_{11}+U_{vv}=0\).
If \(b\ge k^{-12\delta}\), put
\[
 c_0=\frac{b^2H_{11}+H_{vv}}{1+b^2},\qquad U=H-c_0I.
\]
Relative to \(e_1,v,\{e_1,v\}^{\perp}\), a symmetric completion is
\[
 D=\begin{pmatrix}
 -u_{12}/b&bu_{11}&-u_{2E}/b\\
 bu_{11}&bu_{12}&bu_{1E}\\
 -u_{E2}/b&bu_{E1}&0
 \end{pmatrix}.
\]
Here the \(u\)'s are the corresponding blocks of \(U\).
Then \(H-U\ge\tfrac12k^{-2\delta}I\) and
\(|D|\le Ck^{16\delta}\).

For \(0\le b<k^{-12\delta}\), set \(K=k^{14\delta}\) and
\[
 U=-K(I-vv^{\mathsf t})+b^2Kvv^{\mathsf t},\qquad
 D=-bK(e_1v^{\mathsf t}+ve_1^{\mathsf t}).
\]
The identity holds also at \(b=0\). In
\(H-U-k^{-3\delta}I\), the \(v^\perp\) block is at least
\(KI/2\), while the scalar \(v\) block is at least
\[
 k^{-2\delta}-k^{-20\delta}-k^{-10\delta}-k^{-3\delta}.
\]
Its Schur complement subtracts at most \(C|H|^2/K\le Ck^{-6\delta}\).
All negative terms are smaller than \(k^{-2\delta}\) for large \(k\).
The size bounds follow from the explicit matrices.
\end{proof}

Only the Hermitian size of a complex direction is needed to solve
the subsequent polynomial equations. To see this, let \(\mathcal P_m\)
be the homogeneous complex polynomials of degree \(m\), and suppose
\(c\le|q|\le C\). Choose \(|q_i|\ge c/2\), write
\(q\cdot\partial=q_i\partial_i+D'\), and let \(I_i\) integrate
in \(X_i\) with zero constant. The operator
\(\mathcal K=D'I_i/q_i\) lowers the degree in the other variables,
so \(\mathcal K^{m+1}=0\) on \(\mathcal P_m\). Thus
\begin{equation}\label{eq:polyinverse}
 R_{q,m}=\frac{I_i}{q_i}\sum_{r=0}^m(-\mathcal K)^r,\qquad
 (q\cdot\partial)R_{q,m}=I .
\end{equation}
Its coefficient norm is bounded in terms of \(m,c,C\).
There is no division by the possibly small bilinear square \(q\cdot q\).

\subsection{Four packets with arbitrary finite accuracy}

For packets centered at \(z\), define their normalized first jets
in that center's fixed normal coordinates by
\[
 {\bf j}_z w(x)=e^{-kh(x)}
       \left(w(x),\frac{\partial_\xi w(x)}{kQ_z}\right),
 \qquad Q_z=\sqrt{1+|\nabla h(z)|^2}.
\]
The real and imaginary columns of these jets form a real matrix.
Its least spanning singular value means the square root of the
least eigenvalue of the matrix times its transpose.

\begin{proposition}\label{prop:packets}
Given an integer \(J\ge0\) and \(M>0\), finite construction orders
can be fixed so that, for every sufficiently large \(k\), every
center \(z\) with \(t(z)\ge1/4\) has four smooth complex packets.
They have support radius at most \(Ck^{-1/3-100\delta}\), and
\begin{align}
 |D_x^j w|&\le C_j k^{2j}e^{kh}\qquad(j\ge0),\label{eq:packetder}\\
 |D_x^j(L+k^2)w|&\le Ck^{-M}e^{kh}
                    \qquad(0\le j\le J).\label{eq:packetres}
\end{align}
The family spans real first jets with singular value at least
\(ck^{-3/4}\) at \(z\), and at least \(k^{-1}\) at distance
at most \(k^{-10}\) from \(z\). The support, derivative and
spanning exponents do not depend on \(J,M\). Constants and
frequency thresholds may depend on the fixed pair and chosen orders.
\end{proposition}

\begin{proof}
The envelope bounds give \(Q\le k^{3\delta}\) and
\(|H|\le k^{4\delta}\) for large \(k\).
If \(p\ne0\), take \(e_1=p/|p|\); choose a favorable test \(e_2\)
in \eqref{eq:envelope}. At \(p=0\), choose \(e_2\) first and
then any perpendicular \(e_1\). Complete an orthonormal basis and put
\[
 \eta_{\rm v}=k^{-9\delta},\quad v_2=e_2,\quad
 v_3=\frac{e_2+\eta_{\rm v}e_3}{\sqrt{1+\eta_{\rm v}^2}},\quad
 v_4=\frac{e_2+\eta_{\rm v}e_4}{\sqrt{1+\eta_{\rm v}^2}}.
\]
Use the four pairs
\[
 (p',v)=(p,v_2),(p,v_3),(p,v_4),
             (p+Qk^{-3/4}e_1,v_2).
\]
For each set \(Q'=\sqrt{1+|p'|^2}\), \(P=p'+iQ'v\),
\(b=|p'|/Q'\). Then \(P\cdot P=-1\) and
\(1\le Q'/Q\le1+k^{-3/4}\).
The change in the favorable Hessian test is at most
\(Ck^{-5\delta}+Ck^{-3/4+4\delta}\).
For the second term, differentiate \(s^2/(1+s^2)\) as
\(s=|p|\) increases by \(Qk^{-3/4}\).
These errors are smaller than the original \(ck^{-\delta}\)
margin. Lemma~\ref{lem:hessian} therefore gives \(U,D\) with
\((U+iD)P=0\) for each pair.

Each packet uses one frozen normal chart
\(\Psi_{z,E}(\xi)=\exp_z(\sum_i\xi_iE_i)\).
For a fixed smooth input metric, all fixed-order derivatives of
these charts and their inverses are uniformly bounded over centers
and unit frames on sufficiently small common balls. This follows
by solving the smooth geodesic equation on a common bounded set
and differentiating it in initial conditions, then using compactness
of the frame bundle. No spatial derivative of the frame-selection
rule is taken.

Introduce the scaled variables
\begin{equation}\label{eq:scaling}
 r_0=k^{-100\delta},\quad X=\xi/r_0,\quad
 \kappa=kr_0Q,\quad
 h_{\rm sc}(X)=\frac{h(\Psi_{z,E}(r_0X))-h(z)}{r_0Q}.
\end{equation}
For \(q\ge2\), the chain rule gives
\[
 |D_X^q h_{\rm sc}|\le C_q Q^{-1}k^{\delta(101-99q)}.
\]
These are bounded for every fixed order. If
\(L=g^{ij}(\xi)\partial_{\xi_i\xi_j}
       +\beta^i(\xi)\partial_{\xi_i}\), multiplying \(L+k^2\)
by \(r_0^2\) gives
\[
 \mathcal L+\kappa^2/Q^2,\qquad
 \mathcal L=g^{ij}(r_0X)\partial_{X_iX_j}
                   +r_0\beta^i(r_0X)\partial_{X_i}.
\]
Its fixed-order coefficient derivatives are bounded, with
\(g^{ij}(0)=\delta_{ij}\) and \(\partial g^{ij}(0)=0\).

For an integer \(T\ge3\), construct a polynomial phase with
\[
 \Phi(0)=0,\quad \partial\Phi(0)=P/Q,\quad
 \partial^2\Phi(0)=r_0(U+iD)/Q
\]
so that
\(E_\Phi=g^{ij}(r_0X)\partial_i\Phi\,\partial_j\Phi+Q^{-2}\)
vanishes to order \(T\) at zero. The constant and linear equations
are already satisfied. At homogeneous degree \(m\ge2\), the new
term \(\Phi_{m+1}\) enters as
\(2(P/Q)\cdot\partial\Phi_{m+1}\); everything else is known.
Its direction has Hermitian norm between \(1\) and \(3\), so
\eqref{eq:polyinverse} solves every equation.
For fixed \(T\) all coefficients are uniformly bounded.

Put
\[
 \mathcal T a=2g^{ij}(r_0X)(\partial_i\Phi)(\partial_j a)
                          +(\mathcal L\Phi)a.
\]
For a finite integer \(s\ge0\), solve by homogeneous degrees for
polynomials \(a_0,\ldots,a_s\), with \(a_0(0)=1\), \(a_j(0)=0\)
for \(j\ge1\), such that
\[
 \mathcal T a_0,\qquad
 \mathcal T a_j+\mathcal L a_{j-1}\quad(1\le j\le s)
\]
vanish to order \(T\).
Again the new degree-\((m+1)\) term enters via
\(2(P/Q)\cdot\partial\). Complete all of \(a_{j-1}\) before
starting \(a_j\), so the second derivatives in its right-hand side
are already determined. Every coefficient is bounded for fixed \(T,s\).

Choose a real smooth cutoff \(\omega\), equal to one on the unit
ball and supported in the radius-two ball, and define
\begin{equation}\label{eq:packet}
 a=\sum_{j=0}^s\kappa^{-j}a_j,\qquad
 w(\Psi_{z,E}(r_0X))=
 e^{kh(z)}e^{\kappa\Phi(X)}a(X)\omega(k^{1/3}X).
\end{equation}
It extends smoothly by zero and has the stated support radius.
On its support \(a=1+O(k^{-1/3})\), hence \(|a|\ge1/2\).

The real phase lies below the envelope to quadratic order:
\begin{equation}\label{eq:localization}
 \operatorname{Re}\Phi-h_{\rm sc}
 \le k^{-3/4}|X|-c k^{-106\delta}|X|^2+C|X|^3 .
\end{equation}
The linear term is zero for the first three packets. On
\(|X|\le2k^{-1/3}\), the cubic term is absorbed into half the
quadratic deficit. The remaining maximum is
\(O(k^{-3/2+106\delta})\); since
\(k^{1-100\delta}\le\kappa\le k\), the exponential in
\eqref{eq:packet} is at most \(2e^{kh}\).
On the annulus \(\tfrac14k^{-1/3}\le|X|\le2k^{-1/3}\)
it is at most
\begin{equation}\label{eq:tail}
 e^{-2k^{1/4}}e^{kh}.
\end{equation}
Indeed, the negative exponent has size at least
\(ck^{1/3-206\delta}\), and the positive linear exponent is
at most \(Ck^{-1/12}\).

Direct differentiation before the cutoff gives
\[
\begin{split}
 e^{-\kappa\Phi}(\mathcal L+\kappa^2/Q^2)(e^{\kappa\Phi}a)
 ={}&\kappa^2E_\Phi a+\kappa\mathcal T a_0\\
 &+\sum_{j=1}^s\kappa^{1-j}
       (\mathcal T a_j+\mathcal L a_{j-1})
       +\kappa^{-s}\mathcal L a_s .
\end{split}
\]
Taylor remainders and their \(q\le J\le T\) derivatives cost
\(|X|^{T-q}\). In ordinary coordinates a derivative of such a
remainder costs at most \(k^{1/3+100\delta}\);
derivatives of the exponential, amplitude and cutoff cost at
most \(Ck^2\) each. Undoing the operator scaling gives the bound
\begin{equation}\label{eq:accuracy}
 Ck^{200\delta+2J}
 \left(k^{2-T/3}+k^{1-T/3}+k^{-(1-100\delta)s}\right)e^{kh}.
\end{equation}
All cutoff commutators have polynomial factors times \eqref{eq:tail}.
Choose, before \(k\),
\[
 T>3(M+200\delta+2J+3),\qquad
 s>\frac{M+200\delta+2J+1}{1-100\delta}.
\]
Increasing the frequency threshold then proves \eqref{eq:packetres}.
The same calculation without the operator factor proves
\eqref{eq:packetder}, with an exponent independent of \(T,s\).

At the center,
\({\bf j}_z w=(1,P/Q)+O(\kappa^{-1})\).
The real leading columns contain
\((1,(|p|/Q)e_1)\) and
\((1,(|p|/Q+k^{-3/4})e_1)\); they span the value--\(e_1\)
plane with singular value at least \(ck^{-3/4}\).
The imaginary columns span \(e_1^\perp\) with singular value
at least \(ck^{-9\delta}\).
The amplitude error is \(O(k^{-1+100\delta})=o(k^{-3/4})\).
Through two derivatives \eqref{eq:packetder} also gives
\(|D_x{\bf j}_z w|\le Ck^3\). Moving by \(k^{-10}\)
changes the matrix by only \(O(k^{-7})\).
Fixed coordinate changes have bounded inverses, so the nearby
lower bound is \(k^{-1}\) for large \(k\).
Every step includes \(p=0\).
\end{proof}

\subsection{Estimates on wavelength balls}

The next estimates control all packets that meet a small ball,
including packets centered outside the region where signs will be
measured. This is needed because a value-normalized Gaussian sum
must have bounded local derivative moments.

\begin{lemma}\label{lem:scaled}
Fix the packet orders. Let \(\mathcal W\) contain all four packets
at each of finitely many centers, with
\(|\mathcal W|\le Ck^{K_1}\) for fixed \(K_1\).
Suppose every \(x\in V\) has a center within \(C_0k^{-10}\).
Define, in the fixed chart \(\Omega\),
\[
 R_x=(kQ_x)^{-1},\quad
 S_x^2=\sum_{w\in\mathcal W}|w(x)|^2,\quad
 w_x(Y)=w(x+R_xY),\quad |Y|\le2.
\]
For packets whose support meets this ball, let \(X_z\) denote
their own centered scaled coordinates. Call them non-tail when
\(|X_z(x)|<\tfrac12k^{-1/3}\), and tail otherwise.
For every fixed \(j\ge0\),
\begin{align}
 S_x&\ge ce^{kh(x)},\label{eq:variance}\\
 \sup_{|Y|\le2}|D_Y^jw_x(Y)|
 &\le C_j|w(x)| &&\text{for non-tails},\label{eq:nontail}\\
 \sup_{|Y|\le2}|D_Y^jw_x(Y)|
 &\le C_j e^{-k^{1/4}}e^{kh(x)}
                         &&\text{for tails}.\label{eq:abstail}
\end{align}
Every non-tail has nonzero center value and cutoff identically
one on the ball, and
\begin{equation}\label{eq:imaginary}
 \left|\operatorname{Im}\frac{\dd_Yw_x}{w_x}(0)\right|\ge c.
\end{equation}
Consequently,
\begin{equation}\label{eq:squares}
 S_x^{-2}\sum_w\sup_{|Y|\le2}|D_Y^jw_x|^2\le C_j,\qquad
 S_x^{-2}\sum_{\text{non-tail }w}|w(x)|^2=1-o(1).
\end{equation}
Constants and thresholds may depend on the fixed pair, orders
and \(K_1,C_0\), but not on \(k,x\).
\end{lemma}

\begin{proof}
The envelope gives \(|D Q|\le Ck^{4\delta}\).
Any intersecting packet center has
\[
 |z-x|\le Ck^{-1/3-100\delta}+C/k,\qquad
 |Q_z-Q_x|\le Ck^{-1/3-96\delta}+Ck^{-1+4\delta}=o(1).
\]
All such centers and balls lie in a common compact subchart for
large \(k\). Since \(Q\ge1\), the ratios \(Q_z/Q_x\) and the
corresponding ratios throughout the ball are bounded above and
below. From \(|\dd h|\le C Q\),
\[
 k|h(x+R_xY)-h(x)|\le C \quad(|Y|\le2).
\]
The nearest center has \(|X_z(x)|=O(k^{-10+100\delta})\).
Its center value and nearby jet estimates yield
\(|w(x)|\ge ce^{kh(x)}\) for one of its packets, proving
\eqref{eq:variance}.

With \(\kappa_z=kr_0Q_z\), fixed normal-chart bounds give
\[
 |D_Y^qX_z|\le C_qr_0^{-1}R_x^q,\qquad
 \kappa_z|D_Y^qX_z|\le
 C_q(Q_z/Q_x)(kQ_x)^{1-q}\le C_q\quad(q\ge1).
\]
For \(q=1\),
\(\kappa_z D_YX_z=(Q_z/Q_x)D_x\xi_z\) has a uniformly
bounded inverse. The displacement of \(X_z\) on the ball is
\(O(k^{-1+100\delta}/Q_x)=o(k^{-1/3})\).

For a tail, any point where the cutoff or its derivatives
contribute satisfies
\(\tfrac14k^{-1/3}\le|X_z|\le2k^{-1/3}\).
Use \eqref{eq:tail}. Each fixed positive \(Y\)-derivative of
\(\kappa_z\Phi(X_z)\) is bounded by the preceding estimates;
products with \(m\) derivatives of \(X_z\) leave
\(\kappa_z^{1-m}\le1\).
Amplitude derivatives are bounded, and
\(k^{1/3}/\kappa_z\le k^{-2/3+100\delta}\le1\)
controls cutoff derivatives. The exponential estimate therefore
gives \eqref{eq:abstail}, even when \(w(x)=0\).

For a non-tail the cutoff stays one on the ball and the amplitude
is bounded away from zero. The same derivative bounds control
the logarithmic derivative and all higher derivatives relative
to \(|w_x|\). Integrating its real logarithmic derivative along
straight segments gives \(|w_x(Y)|\le C|w(x)|\), proving
\eqref{eq:nontail}.
Moreover,
\[
 \operatorname{Im}\partial_X\Phi(X_z(x))
       =(Q'/Q_z)v+O(k^{-1/3}).
\]
Its leading norm is at least one, so the inverse bound on
\(\kappa_zD_YX_z\) gives a positive imaginary \(Y\)-gradient.
The amplitude's logarithmic derivative is only
\(O(k^{-1+100\delta}/Q_x)=o(1)\), which proves
\eqref{eq:imaginary}. Squaring the tail bounds and using
\eqref{eq:variance} bounds their total normalized contribution
by \(C_j k^{K_1}e^{-2k^{1/4}}=o(1)\).
The non-tail contributions are controlled by their center
values. The definition of \(S_x\) then proves \eqref{eq:squares}.
\end{proof}
